\documentclass[11pt]{article}

\usepackage[T1]{fontenc}
\usepackage[utf8]{inputenc}
\usepackage{amsmath,amssymb,amsthm}
\usepackage[margin=1.15in]{geometry}
\usepackage{microtype}
\usepackage{booktabs}
\usepackage{enumitem}
\usepackage[hidelinks]{hyperref}

% ---- theorem environments -------------------------------------------------
\theoremstyle{plain}
\newtheorem{theorem}{Theorem}
\newtheorem{proposition}[theorem]{Proposition}
\newtheorem{corollary}[theorem]{Corollary}
\newtheorem{lemma}[theorem]{Lemma}
\theoremstyle{remark}
\newtheorem{remark}[theorem]{Remark}

% ---- provenance marks and Lean identifiers --------------------------------
% Every result carries its provenance, because the three kinds differ in what
% they can be held to.  See the legend in Section 1.
\newcommand{\prov}[1]{\textnormal{\footnotesize\textsc{#1}}}
\newcommand{\leanref}[1]{%
  \par\smallskip\noindent{\footnotesize\ttfamily #1}\par}

% ---- notation -------------------------------------------------------------
\newcommand{\ist}{\mathit{ist}}
\newcommand{\Cn}{\operatorname{Cn}}
\newcommand{\Cncirc}{\operatorname{Cn_{circ}}}
\newcommand{\truths}{\operatorname{truths}}
\newcommand{\Mod}{\operatorname{Mod}}
\newcommand{\Lift}{\operatorname{Lift}}
\newcommand{\True}{\mathit{True}}
\newcommand{\CL}{\mathit{CL}}
\newcommand{\binEntropy}{\operatorname{H}_{\mathrm{bin}}}

\title{Harmless and Pointless\\[0.3em]
  \large Consistency, rigidity, and the indeterminacy of context transfer\\
  in McCarthy's formalization of context}
\author{Mateo Estrada}
\date{7 September 2026}

\begin{document}
\maketitle

\begin{abstract}
In \emph{Notes on Formalizing Context} (1993), John McCarthy proposed treating
contexts as first-class objects related by $\ist(c,p)$, observed that this
generates an infinite regress of outer contexts, and asserted without proof that
``the regress is infinite, but \ldots\ it is harmless.'' We prove it, and then
show what the proof costs. The regress is harmless precisely because it is
\emph{rigid}: the transcendence schema does not merely permit the model in which
each context copies the one below, it forces it. That rigidity is fatal to the
case McCarthy identifies as interesting, in which a transcending context relaxes
an assumption of the old; we show the two are formally incompatible rather than
merely different. We then examine the natural repairs. Minimizing change over
valuations is degenerate: the minimum always exists, is unique, and changes
nothing. Taking contexts to be theories rescues the phenomenon---dropping an
assumption can create a conclusion---but the accompanying lifting rule is doubly
underdetermined: circumscription admits incomparable minimal resolutions, and
transfer is \emph{route-sensitive}: a skeptical consequence of the direct route
becomes merely credulous after passing through an intermediary. A final finite
extension proves that the two routed blocks exhaust the actual minimal-block
subtype, imposes a uniform policy on that subtype, and derives Shannon
quantities from an explicit joint law.  Direct selection has entropy $0$ while
routed selection has entropy $\log 2$.  Its separate loss calculation uses a
memoryless, fresh-stage channel model.  The finite-sum mutual information of
its parameterized joint law equals the entropy closed form and has exact
one-use capacity $\log(5/4)$ nats, uniquely at routed-input probability $2/5$;
its finite memoryless product channel is normalized before any coding question
is posed.  All 23 results are machine-checked in Lean 4 and free of
\texttt{sorry}. The development is available at
\url{https://github.com/TheodoreAI/mccarthy-contexts} under an MIT license.
\end{abstract}

% ===========================================================================
\section{The claim that was never proved}

McCarthy's proposal is that a proposition is true \emph{in a context}, written
$\ist(c,p)$, and that contexts are objects one can reason about. The difficulty
he immediately notices is that an assertion of $\ist(c,p)$ is itself made
somewhere---in some outer context---so one has ``something like
$\ist(c_0, \ist(c,p))$'', and so on without end.

Two passages carry the claim we address. In \S1 he writes that ``the regress is
infinite, but we will show that it is harmless.'' The promised demonstration
does not appear. In \S5, ``Transcending Contexts,'' the same regress is given
its operational form: believing $p$ unconditionally is taken to be equivalent to
$\ist(c_0,p)$, and the step may be repeated indefinitely.

McCarthy is explicit that his proposals are ``incomplete and tentative'' and
that he will use formulas that ``aren't always compatible with each other.'' A
consistency claim needs a precise system to attach to, so the work below first
fixes one, proves the claim for it, and is then careful about how far the result
reaches.

\paragraph{Provenance.}
Every result is marked, since the three kinds differ in what they can be held
to. \prov{McCarthy} marks a claim he asserts and we prove. \prov{Analysis} marks
a proposal he makes in passing without working out, where the formulation is
ours and a reader should check it against his text. \prov{New} marks results
that go beyond his paper altogether. Every result is additionally verified by
the Lean kernel; Section~\ref{sec:formalization} says what that does and does
not buy.

% ===========================================================================
\section{The tower, and why it is harmless}

Fix a countable set of propositional atoms and build a cumulative tower of
languages $L_0 \subseteq L_1 \subseteq \cdots$, where $L_{n+1}$ adjoins one new
atom $\ist_n(\varphi)$ for each $\varphi \in L_n$. The transcendence schema
$\Gamma$ is the set of sentences
\[
  \ist_n(\varphi) \leftrightarrow \varphi,
  \qquad n \ge 0,\ \varphi \in L_n .
\]

The decisive feature is that $\ist_n$ is \emph{typed}: it applies only to
formulas of level $n$. No formula can apply $\ist_n$ to itself, so no analogue
of the Liar sentence---and no analogue of Russell's paradox for
$\ist$---can be written down at all. This is Tarski's remedy \cite{tarski}, and
McCarthy's own appeal to reflection principles points at it; what was missing
was carrying the analogy out.

\begin{theorem}[\prov{McCarthy}]\label{thm:one}
$\Gamma$ is satisfiable, hence consistent.
\end{theorem}
\leanref{TranscendenceTower.transcendence\_tower\_satisfiable}

The proof is an increasing chain of classical valuations. $M_0$ assigns truth
values to the base atoms arbitrarily; $M_{n+1}$ keeps every old value and sets
$M_{n+1}(\ist_n(\varphi)) = M_n(\varphi)$. This is well defined because
$\varphi$ ranges only over $L_n$, on which $M_n$ was fixed one stage
earlier: each new atom is read off a valuation settled before the new language
existed. No compactness argument is needed.

\begin{corollary}[\prov{McCarthy}]
For every ordinal $\alpha$ and every assignment to the base atoms, the
recursively induced compositional valuation satisfies the tower through
stage $\alpha$.
\end{corollary}
\leanref{TranscendenceTower.Transfinite.Gamma\_satisfiable\_compositional}

McCarthy raises transfinite iteration and sets it aside. The mathematics extends
without new ideas, but the \emph{formalization} does not: limit stages must be
handled explicitly, and that is the one place the construction could fail.

\begin{lemma}[\prov{McCarthy}]
For a limit ordinal $\lambda$, every formula of $L_\lambda$ already belongs to
some $L_\beta$ with $\beta < \lambda$.
\end{lemma}
\leanref{TranscendenceTower.Transfinite.limit\_stage\_union}

A formula is a finite object: it mentions finitely many levels, all below the
limit. So at a limit stage nothing new is decided---the stage only collects what
earlier stages already settled.

% ===========================================================================
\section{Reifying a level}

McCarthy suggests the formalism ``might be further extended to provide so that
in $c_{-1}$ the whole set of sentences true in $c_0$ is an object
$\truths(c_0)$.'' That is the whole of his text on the matter: one sentence,
with no axioms and no claims attached. The results of this section are therefore
analysis rather than transcription---the formulations are ours, and a reader
assessing faithfulness has only that sentence to compare them against. Naming
the set turns out to be free; claiming it is \emph{complete} does not.

\begin{proposition}[\prov{Analysis}]
The Tarski schema $\True_0(\ulcorner\varphi\urcorner) \leftrightarrow \varphi$
is consistent and conservative over $L_0$: every valuation of the base language
is the reduct of a model of it.
\end{proposition}
\leanref{TranscendenceTower.Reification.T0\_consistent\_and\_conservative}

Tarski's undefinability theorem is sometimes misread as forbidding this. It
forbids a sufficiently strong language from defining \emph{its own} truth
predicate; here $\True_0$ sits one level above the language it describes, which
is the stratified arrangement Tarski proposed as the remedy.

\begin{proposition}[\prov{Analysis}]\label{prop:closure}
The closure axiom $\CL(S)$---``the extension of $\True_0$ is exactly
$S$''---is not conservative. There are a satisfiable such theory and a base
sentence it refutes although the base language leaves it open.
\end{proposition}
\leanref{TranscendenceTower.Reification.closure\_not\_conservative}

The three-part shape of that statement is not decoration; see
Section~\ref{sec:formalization}. The negative information a closure axiom
supplies---\emph{nothing else is true here}---is precisely what definitional
extension cannot deliver, and nonmonotonic machinery of this kind is McCarthy's
own invention \cite{circumscription}: $\CL(S)$ is a circumscription of
$\True_0$.

\begin{theorem}[\prov{Analysis}]
If $S \neq S'$ then $T_0 + \CL(S) + \CL(S')$ is inconsistent.
\end{theorem}
\leanref{TranscendenceTower.Reification.closure\_collide}

Two completeness claims about the same level, each honestly made, cannot both be
held merely because each was made in good faith. This bites a single agent as
much as two: a completeness claim made today and another made tomorrow, about
the same level after it has changed, are formally contradictory. The repair is
to relativize---a completeness claim must carry the state it was computed
against---and relativized claims about distinct states coexist
(\texttt{relativized\_satisfiable}).

% ===========================================================================
\section{Harmless because rigid}

Here the paper goes beyond anything in McCarthy's text. The consistency result
of Section~2 understates what the schema does, and the stronger statement
explains both why the tower is safe and why it is uninteresting.

\begin{theorem}[\prov{New}]\label{thm:categorical}
$\Gamma$ has a unique expansion over each fixed base valuation: any two
contexts extending the same level-0 valuation and satisfying $\Gamma$ are
equal.
\end{theorem}
\leanref{TranscendenceTower.Relaxation.schema\_unique}

The schema does not merely permit the copy model; it forces it. This is the
precise content of the observation that each stage is ``a definitional extension
in disguise'': every new atom arrives with an axiom defining it outright in
older vocabulary. Definitional extensions are conservative and preserve
consistency---which is exactly why the tower is harmless, and exactly why
McCarthy calls the resulting picture ``pointless'', the new sentences being
``more elaborate versions of the old.''

\begin{theorem}[\prov{New}]
No context both satisfies $\Gamma$ and relaxes an assumption of the level below.
\end{theorem}
\leanref{TranscendenceTower.Relaxation.no\_relaxing\_model\_of\_schema}

So the two cases McCarthy distinguishes are not two situations to be handled in
turn. They are formally exclusive. This explains, rather than restates, why a
consistency proof for the tower cannot reach the interesting case: the obstacle
is not that the technique is too weak but that the statement being proved
concerns a schema which relaxation contradicts. The loss is immediate---dropping
a single instance already admits two distinct extensions
(\texttt{relaxed\_not\_definitional}), so no weakened form of definitional
extension survives partial relaxation.

% ===========================================================================
\section{Minimal change is degenerate}

If relaxation leaves several contexts available, the natural repair is to prefer
the one that changes least---circumscription, which is McCarthy's own
instrument. Write $\Delta(N)$ for the set of assumptions at which a transcending
context departs from the one below, and $\Mod(v,A)$ for the contexts obeying the
schema outside a relaxation set $A$.

\begin{theorem}[\prov{New}]\label{thm:bijection}
$\Delta$ is a bijection from $\Mod(v,A)$ onto the subsets of $A$. Relaxing $k$
assumptions leaves exactly $2^k$ contexts, and the relaxed schema constrains
nothing beyond the location of change.
\end{theorem}
\leanref{TranscendenceTower.MinimalChange.delta\_bijection}

Injectivity turns on truth values being \emph{Boolean}: knowing where a context
departs from the base determines what it says there, because only one
alternative exists. That observation makes the next result inevitable, and it
also tells us where to go afterwards.

\begin{theorem}[\prov{New}]\label{thm:collapse}
For every relaxation set, the copy context is the unique $\Delta = \emptyset$
model and hence the $\subseteq$-minimum. Preferring minimal change returns the
unrelaxed tower no matter what was relaxed.
\end{theorem}
\leanref{TranscendenceTower.MinimalChange.copy\_is\_minimum}

Circumscribing the abnormality recovers precisely the case McCarthy calls
pointless. Permitting a context to drop an assumption never forces it to, and
the minimal-change reading always declines. By contrast, \emph{specifying} the
change set determines the context uniquely (\texttt{forced\_unique}). The
content of a logic of relaxing contexts therefore cannot come from a preference
order over these models; it must come from saying which assumptions are dropped.

% ===========================================================================
\section{Contexts as theories}

A valuation has no interior: each atom is independent, and changing one moves
nothing else. So we change what a context \emph{is}. Let a context be a set of
sentences and let ``true in the context'' mean derivable from it. Now dropping
an assumption alters a whole consequence set rather than flipping one bit---and
a minimality principle has something to act on.

\begin{theorem}[\prov{New}]\label{thm:nonmono}
Classical consequence is monotone; circumscriptive consequence is not. There are
$T' \subseteq T$ and $\varphi$ with $\varphi \in \Cncirc(T')$ and
$\varphi \notin \Cncirc(T)$.
\end{theorem}
\leanref{TranscendenceTower.ContextTheories.circ\_nonmonotone}

The witness is the standard default: \emph{birds fly unless abnormal}, plus a
bird. Nothing forces abnormality, so the most normal world has none, and the
bird flies. Assert abnormality and the conclusion is withdrawn---more
information, fewer conclusions. Read in the direction of relaxation,
\emph{removing} the abnormality assumption \emph{creates} the conclusion, which
is the dropped-assumption phenomenon McCarthy wants.

One guard matters: the conclusion is \emph{not} a classical consequence of the
smaller theory (\texttt{default\_reasoning\_works}). Without that check one
could object that circumscription merely inherited what logic already gave.

Note what changed and what did not. The \emph{same} minimality principle that
did nothing over valuations does real work over theories. One variable was
changed, and the outcome reversed.

% ===========================================================================
\section{Lifting, and two failures of determinacy}

McCarthy's \emph{specializes} axiom transfers facts between contexts unless
abnormal for the transfer:
\[
  \mathit{specializes}(c_1,c_2) \wedge \neg \mathit{ab}(p,c_1,c_2)
  \wedge \ist(c_1,p) \rightarrow \ist(c_2,p).
\]
We model a transfer by a \emph{blocking set}---the assertions declared
abnormal---and circumscription as preferring the smallest one. Nothing here
assumes the contexts are stacked, so this covers sibling contexts, not only
tower levels. Unblocked, lifting delivers what it promises
(\texttt{lift\_transfers}); a conflicting target forces a block and the fact
does not survive it (\texttt{lift\_defeasible}). Then two things go wrong.

\begin{theorem}[\prov{New}]\label{thm:extensions}
A source asserting $p$ and $q$ lifted into a target asserting
$\neg(p \wedge q)$ admits two incomparable minimal blocking sets. Minimizing
abnormality does not determine the lifted context.
\end{theorem}
\leanref{TranscendenceTower.Lifting.lift\_multiple\_extensions}

This is the familiar multiple-extensions phenomenon of nonmonotonic reasoning,
and we claim no novelty for it. What the result contributes is its
\emph{location}: the failure occurs inside the rule that transfers facts between
contexts, on the same semantics used throughout, at the point where a theory of
contexts most needs a determinate answer.

\begin{theorem}[\prov{New}]\label{thm:route}
With $c_1 \vdash p$, $c_2 \vdash q$ and $c_3 \vdash \neg(p \wedge q)$: directly,
$p$ is a consequence under every minimal resolution. Through $c_2$, two
incomparable minimal resolutions appear and $p$ survives only one; it is
therefore merely a credulous consequence of the routed transfer.
\end{theorem}
\leanref{TranscendenceTower.Chaining.lifting\_is\_route\_sensitive}

The endpoints do not conflict: $p$ and $\neg(p \wedge q)$ hold together. The
indeterminacy is manufactured entirely by the route, because the first hop
deposits $c_2$'s own commitment into the destination alongside $p$. So
\emph{specializes} composes as a relation between contexts while the transfer it
licenses does not: the composite is not determined by its endpoints, and lifting
cannot be chained freely to reach a distant context.

% ===========================================================================
\section{A finite information calculation}

The formal proof first establishes that the two displayed routed blocks
exhaust the actual subtype of minimal blocks: every such block is either
``drop $p$'' or ``drop $q$''.  They are alternatives, not outcomes of a
probability space.  To make a finite quantitative statement, we therefore
\emph{add} a uniform law on that actual subtype.  Entropy applies only after
imposing this policy; it is not an entropy of nondeterministic consequence or
of the lifting relation by itself.

\begin{theorem}[\prov{New}]\label{thm:information}
The unique direct minimal block carries a singleton (Dirac) law and hence has
entropy $0$.  Under the fair routed policy, routed selection has entropy
$\log 2$ nats (one bit).  For a uniform prior on direct versus routed transfer,
an explicit four-cell route/output joint law has survival marginal $3/4$,
conditional output entropy $(\log 2)/2$, and mutual information
\[
  \binEntropy(3/4) - (\log 2)/2.
\]
These Shannon quantities are finite sums derived from that joint law.  In the
separate abstract loss-only channel, if $0\leq s,t\leq1$, composing
memoryless fresh-stage Markov kernels with survival probabilities $s$ and $t$
gives the valid kernel with survival probability $st$; hence direct then
routed survival is $1/2$, and two fresh fair routed stages have survival
$1/4$.
\end{theorem}
\leanref{TranscendenceTower.InformationTheory.minblock\_routed\_iff;
TranscendenceTower.InformationTheory.direct\_selection\_entropy;
TranscendenceTower.InformationTheory.symmetric\_routed\_selection\_entropy;
TranscendenceTower.InformationTheory.routeOutputJointMass\_is\_distribution;
TranscendenceTower.InformationTheory.uniform\_route\_survival\_probability;
TranscendenceTower.InformationTheory.uniform\_route\_conditional\_entropy;
TranscendenceTower.InformationTheory.uniform\_route\_mutual\_information;
TranscendenceTower.InformationTheory.fresh\_stage\_lossChannel\_compose;
TranscendenceTower.InformationTheory.two\_routed\_fresh\_stages\_survival}

Here ``nats'' records the natural-logarithm convention: $\log 2$ nats is one
bit.  This is intentionally a finite calculation about the proved actual
minimal-block subtype and its added policy.  It makes no coding-theorem claim,
does not turn arbitrary nonmonotonic choices into probabilities, and does not
infer multi-stage behavior from one-stage logical marginals.

\subsection{Capacity and memoryless repeated use}

Let $q\in[0,1]$ be the probability of choosing the routed input, so the direct
input has probability $1-q$.  Since the route alphabet has exactly two
elements, every input probability distribution is uniquely of this form.
Write $J_q$ for the resulting route/output joint law, $p_q$ for its route
prior, and $m_q$ for its output marginal.  The
formal finite-sum mutual information, with zero-mass cells contributing zero,
is proved equal to the entropy closed form:
\[
  I_{\mathrm{joint}}(q)
  =\sum_{r,b:J_q(r,b)\ne0}J_q(r,b)
    \log\frac{J_q(r,b)}{p_q(r)m_q(b)}
  =\binEntropy(q/2)-q\log 2.
\]
The proof evaluates $q=0$ and $q=1$ separately and uses nonzero denominators
only in the interior $0<q<1$.

\begin{theorem}[\prov{New}]\label{thm:capacity}
For every $q\in[0,1]$,
\[
  I_{\mathrm{joint}}(q)\leq\log(5/4),
  \qquad
  I_{\mathrm{joint}}(q)=\log(5/4)\ \Longleftrightarrow\ q=2/5.
\]
In particular, $I_{\mathrm{joint}}(2/5)=\log(5/4)$, so the channel capacity is
$\log(5/4)$ nats and the optimizing routed-input probability is unique.
For every input word $x:\mathrm{Fin}\ n\to\{\mathsf{direct},\mathsf{viaMid}\}$,
the repeated-use channel
\[
  K_n(x,y)=\prod_{i:\mathrm{Fin}\ n}K(x_i,y_i)
\]
is a nonnegative normalized mass function over output words
$y:\mathrm{Fin}\ n\to\{\mathsf{false},\mathsf{true}\}$.  The cases $n=0$ and
$n=1$ are respectively the unit mass on the empty word and the original
one-use channel.
\end{theorem}
\leanref{TranscendenceTower.InformationTheory.routePrior\_eq\_routedInputPrior;
TranscendenceTower.InformationTheory.oneUseJointMutualInformation\_eq\_closed\_form;
TranscendenceTower.InformationTheory.oneUseJointMutualInformation\_at\_two\_fifths;
TranscendenceTower.InformationTheory.oneUseJointMutualInformation\_capacity;
TranscendenceTower.InformationTheory.oneUseJointMutualInformation\_eq\_capacity\_iff;
TranscendenceTower.InformationTheory.blockChannel\_is\_distribution;
TranscendenceTower.InformationTheory.blockChannel\_zero;
TranscendenceTower.InformationTheory.blockChannel\_one}

The product normalization establishes precisely the memoryless finite channel
required before introducing codes.  It neither constructs an encoder or
decoder nor proves an achievability, converse, or asymptotic coding theorem.

% ===========================================================================
\section{What the three failures have in common}

Minimizing abnormality fails to deliver a context three times, in three
different ways, and the pattern is the argument of this paper.

\begin{center}
\begin{tabular}{@{}lll@{}}
\toprule
Setting & Minima & Failure \\
\midrule
Valuations & unique & always the trivial one---no information \\
Theories, one transfer & plural & incomparable---no choice made \\
Theories, composed transfers & route-relative & not determined by the endpoints \\
\bottomrule
\end{tabular}
\end{center}

\medskip

The first two are complementary rather than repetitive: over valuations
minimization is uninformative because the minimum is trivial, over theories
because the minima are several. The third is worse than either, since it defeats
the natural assumption that a relation between contexts can be followed
transitively.

None of this shows that no logic of relaxing contexts exists. It locates
precisely what such a logic must supply and cannot obtain from minimality alone:
a priority among defaults, or an explicit record of which assumptions were
dropped and by which route. That is a constraint on the programme, not a defect
of these examples.

% ===========================================================================
\section{On the formalization}\label{sec:formalization}

All results are Lean 4 (toolchain \texttt{v4.33.1}, Mathlib \texttt{0df444a3}),
free of \texttt{sorry}, \texttt{admit} and \texttt{native\_decide}, with axiom
use confined to \texttt{propext}, \texttt{Classical.choice} and
\texttt{Quot.sound}. Several depend on no axioms at all.

The development is public at
\url{https://github.com/TheodoreAI/mccarthy-contexts} under an MIT license, with
the toolchain and Mathlib revision pinned to the environment above.
\texttt{lake build} checks every result; \texttt{scripts/verify.sh} runs the
whole gate---build, a scan for unproved holes, and the axiom audit---and exits
non-zero if any part fails, so the claims of this section can be rechecked
rather than taken on trust.

Machine checking prevents gaps in proofs. It does not prevent a proof of the
wrong thing, and the two failures worth reporting are of that kind.

\paragraph{A vacuous theorem that compiled.}
An early statement of Proposition~\ref{prop:closure} read: \emph{every model of
$T_0 + \CL(\emptyset)$ refutes $\varphi$}. It compiled, contained no
\texttt{sorry}, and was worthless. In any model of $T_0$ the extension of
$\True_0$ contains every tautology, so it is never empty and the theory has no
models at all; a refutation claim about a theory with no models establishes
nothing. The published form therefore asserts \emph{satisfiability} as an
explicit conjunct.

This motivated a standing check applied to every result here: for any statement
quantifying over models, prove the model class non-empty; for any hypothesis
defined in-file, exhibit a witness. The transfinite development defines its own
notion of limit ordinal, so it also proves that $\omega$ satisfies it---otherwise
the limit-stage lemma would be quietly empty. Its satisfiability theorem also
names the recursively defined valuation in the statement itself; quantifying
over an arbitrary function from formulas to Booleans would admit the
constant-true function without enforcing the semantics of any connective.

\paragraph{Naming what is not proved.}
The second discipline is scope. Theorem~\ref{thm:one} is a statement about a
\emph{stratified propositional} system in which each context copies the one
below. It is not a proof that McCarthy's formalism is consistent, which cannot
currently be evaluated because it was never made precise. Extending the
language---in particular by quantifying over contexts as first-class
terms---would not falsify Theorem~\ref{thm:one}; it would produce a different
system to which the theorem does not apply. Since stratification is what the
proof rests on, that extension is also the one most likely to break it, and it
should be attempted deliberately rather than by accident.

\paragraph{Disclosure: AI assistance.}
The Lean development and the text of this paper were produced with Claude
(Anthropic, Opus~5), in a session directed by the author, who set the targets,
chose the formulations and the scope of each claim, and decided which lines of
enquiry to pursue. The standing of the results does not rest on that assistance.
Every theorem is checked by the Lean kernel, which is independent of the process
that produced the proof: a proof term either type-checks against the stated type
or it does not, and the axiom audit is machine-generated rather than asserted. A
reader who distrusts the provenance entirely can recheck the development and
lose nothing. What machine checking does not supply is assurance that the
\emph{intended} statement was proved, and that is where assistance of this kind
can go wrong silently---the vacuous formulation described above arose in exactly
this way, and was caught by the discipline of this section rather than by
reading the proof.

% ===========================================================================
\section{Open}

\begin{itemize}[leftmargin=1.2em]
  \item \textbf{Priorities among defaults.} The direct response to Section~7:
    does a priority ordering restore a unique lifted context, and under exactly
    what condition?
  \item \textbf{Quantification over contexts.} First-class context terms, faced
    together with the self-reference risk described in
    Section~\ref{sec:formalization}.
  \item \textbf{The theory-level tower.} The results of Section~6 concern one
    step. Whether the nonmonotonic behaviour survives iteration, and what
    happens at limit stages, is untouched.
\end{itemize}

% ===========================================================================
\begin{thebibliography}{9}

\bibitem{notes}
J.~McCarthy. \emph{Notes on Formalizing Context.} Stanford, 1993
(\texttt{context-2.tex}, created winter 1991).

\bibitem{circumscription}
J.~McCarthy. Circumscription---a form of non-monotonic reasoning.
\emph{Artificial Intelligence} 13 (1980), 27--39.

\bibitem{generality}
J.~McCarthy. Generality in Artificial Intelligence. \emph{CACM} 30:12 (1987),
1030--1035.

\bibitem{tarski}
A.~Tarski. Der Wahrheitsbegriff in den formalisierten Sprachen.
\emph{Studia Philosophica} 1 (1936).

\bibitem{buvac93}
S.~Buva\v{c} and I.~A.~Mason. Propositional logic of context. \emph{AAAI-93},
412--419.

\bibitem{buvac96}
S.~Buva\v{c}. Quantificational logic of context. \emph{AAAI-96}.

\bibitem{feferman}
S.~Feferman. Transfinite recursive progressions of axiomatic theories.
\emph{JSL} 27:3 (1962), 259--316.

\end{thebibliography}

\medskip
\noindent{\footnotesize Buva\v{c} and Mason \cite{buvac93} prove soundness,
completeness and decidability for a considerably richer propositional logic of
context than the fragment treated here, with a single modality $\ist(k,\varphi)$
and per-context vocabularies. Their result is strictly more general on that
axis; the development above is correspondingly more elementary, needing nothing
beyond direct limits of classical valuations.}

\end{document}
